% Distribución íntegra por residencia del propietario cosmológico.
\section{Involution des feuillets dans une région piégée et représentation causale de l'intérieur d'un trou noir}

Dans la région extérieure vide d'un corps compact, l'équation de Cartan impose
une torsion nulle lorsqu'il n'existe pas de courant de spin. La connexion se réduit à
\(\mathring\omega(e)\) et la métrique satisfait le secteur d'Einstein. Les orbites,
l'effet de lentille gravitationnelle, le décalage vers le rouge et la propagation extérieure restent
régis par la géométrie relativiste dans le domaine vérifié. HMT--MD situe son
apport dans le sens de la frontière et dans la continuité de l'état qu'une
carte extérieure cesse de publier.

Un horizon est, dans cette lecture, une frontière de publication. L'état
enrichi ne se termine pas là où l'observateur extérieur perd une coordonnée
inversible. La relation entre feuillet, orientation et lecture change. Une section qui
franchit la frontière peut cesser d'avoir un représentant dans la carte extérieure et
continuer sur un feuillet conjugué. L'expression \emph{inversion de feuillet} désigne ce
changement : extérieur et intérieur ne sont pas deux substances, mais deux domaines d'une
présentation globale dont le raccordement conserve la route et la mémoire.

Le ruban de Möbius offre le modèle topologique minimal. Un tour transporte la
section vers son feuillet conjugué ; deux tours rétablissent l'orientation. L'observateur
local peut publier une seule branche, tandis que l'état global conserve les deux. Dans
une réalisation relativiste, la même opération doit agir sur les cadres
causaux, la charge et l'orientation temporelle sans détruire la norme. L'horizon
devient ainsi le lieu où une orientation visible devient mémoire de
retour.

L'inversion de feuillet contient une explication possible du paradoxe de l'
information. Il convient de distinguer trois opérations. L'évolution fondamentale peut
être réversible :
\[
 U_{\rm tot}^{-1}U_{\rm tot}=I.
\]
La carte extérieure peut identifier des états différents :
\[
 q_{\rm ext}(x)=q_{\rm ext}(y),
 \qquad x\neq y.
\]
Et l'appareil extérieur peut effacer son registre pour se préparer de nouveau. La première
opération conserve l'information ; la seconde en perd l'accès ; la troisième a
un coût thermodynamique. La disparition apparente derrière l'horizon peut appartenir
à la fibre de \(q_{\rm ext}\) sans exiger une destruction dans \(U_{\rm tot}\).

La lecture par feuillets permet d'écrire un modèle minimal de cette information. Les
canaux de l'angle et du contre-angle sont
\[
 q_\pm
 =
 \exp\!\left(-A_{\rm rad}\mp \Cstar_{\rm rad}\right).
\]
Après normalisation,
\[
 p_\pm=\frac{q_\pm}{q_++q_-},
\qquad
 \rho_{\rm hoja}
 =
 p_+|+\rangle\langle+|
 +
 p_-|-\rangle\langle-|.
\]
L'entropie réduite est
\[
 S_{\rm hoja}
 =
 -\sum_{\pm}p_\pm\log p_\pm
 =
 \log\!\left(2\cosh\Cstar_{\rm rad}\right)
 -
 \Cstar_{\rm rad}\tanh\Cstar_{\rm rad}.
\]
La phase commune \(A_{\rm rad}\) s'annule : le mélange visible dépend de la séparation
entre feuillets. Si la factorisation avancée--retardée existe,
\[
 |\Psi\rangle
 =
 \sqrt{p_+}\,|+\rangle|\mathrm{adv}\rangle
 +
 \eexp^{i\chi_0}
 \sqrt{p_-}\,|-\rangle|\mathrm{ret}\rangle,
\]
cette entropie est l'intrication de deux orientations d'un état pur. Un
horizon peut alors modifier l'accessibilité des branches sans transformer l'
état global en un mélange fondamental.

La thermodynamique exige de garder la température séparée de la durée. La
constante de Boltzmann est un morphisme entre la droite thermique et la droite énergétique,
\[
 k_B:L_\Theta\longrightarrow L_E.
\]
Pour l'horloge à \(108\) états,
\[
 \omega_0=\frac{2\pi}{108\,t_0},
\qquad
 k_B\Theta_{\rm clk}\log3
 =
 \hbar_{\rm int}\omega_0.
\]
L'information d'une cellule ternaire, l'échelle thermique et le quantum d'action sont
ainsi reliés sans confondre \(\log3\) avec l'entropie topologique
\(\log\phiHMT\).

La limite de Landauer s'applique lorsque l'observateur efface :
\[
 Q\geq k_B\Theta\,\Delta S.
\]
Un cycle global réversible qui conserve le registre peut avoir une variation logique
nulle. Un appareil qui réinitialise une mémoire distinguable doit dissiper. Ce que l'on
appelle le Landauer nul est le régime sans effacement, non une machine thermique gratuite.
Dans un trou noir, cette distinction demande si le rayonnement redistribue une
mémoire conservée, si l'extérieur la maintient inaccessible pendant un certain temps ou s'il
existe réellement un canal non unitaire.

L'interprétation devient une théorie physique lorsqu'elle construit une géométrie.
L'extension doit retrouver l'extérieur relativiste, produire une surface
piégée et un horizon régulier, transporter l'état enrichi à travers la
frontière et assigner un recensement d'états à son aire. La continuation doit en outre
produire un signal distinguable : une phase des modes quasi normaux, une structure
de corrélation dans le rayonnement, une polarisation ou un écho dont la valeur n'est pas choisie
après observation. Si l'inversion ne change que les mots \emph{intérieur} et
\emph{extérieur}, elle ne réalise pas encore le trou noir HMT. Si elle satisfait ces
conditions, le problème d'information devient un problème contrôlable de
changement de carte.


